

\section{L'absence de transfert automatique : La critique de la justification utilitariste}
space{0.3cm}
oindent

La question de la place des langues anciennes dans les curriculums scolaires contemporains ne cesse
de soulever des débats passionnés, cristallisant des enjeux qui dépassent largement le cadre de la
simple pédagogie linguistique. Au cœur de cette controverse réside une affirmation tenace, répétée à
l'envi par les défenseurs des humanités classiques depuis la fin du XIXe siècle : l'étude du latin
et du grec ancien constituerait une « gymnastique de l'esprit » irremplaçable, capable de forger une
rigueur intellectuelle transférable aux domaines les plus éloignés, et singulièrement aux
mathématiques.\cite{II-1-2-ref1} Cette section du rapport, intitulée « L'absence de transfert
automatique : faire du grec 'pour la rigueur' ne rend pas rigoureux en maths », se propose de
déconstruire méthodiquement ce postulat à la lumière des sciences cognitives, de la sociologie de
l'éducation et de l'histoire des disciplines scolaires.

Historiquement, la maîtrise des langues classiques n'avait nul besoin de justification utilitariste
indirecte : elle était la clé d'accès indispensable aux savoirs théologiques, juridiques, médicaux
et scientifiques. Le latin était la langue de l'université, de l'Église et de la diplomatie ; le
grec, celle des Évangiles et de la science fondamentale. Cependant, la sécularisation des sociétés
occidentales et la montée en puissance des sciences expérimentales et des langues vivantes au XIXe
siècle ont progressivement érodé cette nécessité pragmatique.\cite{II-1-2-ref3} Face à ce recul, la
légitimation de l'enseignement classique a opéré une mutation stratégique majeure, passant d'une
utilité de contenu (lire les textes) à une utilité de formation (former l'esprit). C'est ainsi
qu'est née la doctrine de la « discipline formelle », postulant que l'effort consenti pour maîtriser
la morphologie complexe du grec ancien musclerait le cerveau de manière générique, conférant à
l'élève une capacité d'analyse logique applicable universellement.\cite{II-1-2-ref4}

Pourtant, l'examen rigoureux de la littérature scientifique accumulée depuis plus d'un siècle
contredit formellement cette intuition. De la réfutation initiale par Thorndike et Woodworth en 1901
aux méta-analyses contemporaines sur le transfert cognitif et l'entraînement cérébral, les preuves
convergent vers une conclusion dérangeante pour l'apologétique classique traditionnelle : le
transfert automatique de compétences entre la grammaire et les mathématiques est une
illusion.\cite{II-1-2-ref6} Ce rapport s'attachera à démontrer que la corrélation indéniable
observée entre l'étude du grec et la réussite en mathématiques ne relève pas d'une causalité
cognitive directe, mais de mécanismes sociologiques de sélection, de distinction et de reproduction
des élites.\cite{II-1-2-ref9}

Nous analyserons d'abord l'effondrement épistémologique de la théorie de la discipline formelle et
l'émergence des concepts de transfert proche et lointain. Nous examinerons ensuite les données
empiriques modernes qui réfutent l'hypothèse d'un gain cognitif généralisé attribuable au latin ou
au grec. Nous plongerons également au cœur de la dissonance cognitive entre la « rigueur »
grammaticale et la « rigueur » mathématique, en nous appuyant sur la linguistique computationnelle
et les neurosciences. Enfin, nous replacerons cet argumentaire dans son contexte historique et
politique, révélant comment il fut construit pour sauver une discipline en quête de légitimité lors
des grandes réformes éducatives de la IIIe République.\cite{II-1-2-ref11}

Pour comprendre pourquoi l'équation « Grec ancien \= Rigueur mathématique » est fausse, il est
impératif de remonter aux fondements de la psychologie de l'éducation qui ont, dès le début du XXe
siècle, invalidé le modèle du « muscle mental ».



\subsection{De la psychologie des facultés à la révolution de thorndike (1901)}

La pédagogie traditionnelle, héritée de la scolastique et revivifiée au XIXe siècle, reposait sur la
« psychologie des facultés ». Selon cette conception, l'esprit humain serait composé de modules
distincts et généraux tels que la Mémoire, le Raisonnement, l'Attention ou la
Volonté.\cite{II-1-2-ref5} Dans cette optique, le contenu de l'enseignement importait moins que la
difficulté de l'exercice : apprendre par cœur des déclinaisons latines ou des paradigmes de verbes
grecs n'avait pas pour but principal la communication ou la compréhension littéraire, mais le «
durcissement » de la faculté de mémoire et de discipline. Comme un haltère renforce le biceps pour
n'importe quelle tâche physique, le grec était censé renforcer la « faculté de raisonnement » pour
n'importe quelle tâche intellectuelle.\cite{II-1-2-ref1}

Ce paradigme a été fracassé par les travaux pionniers d'Edward Thorndike et Robert Woodworth en
1901, dans ce qui reste l'une des expériences les plus célèbres de l'histoire de la psychologie. Ils
ont testé l'hypothèse selon laquelle l'entraînement intensif d'une fonction mentale spécifique (par
exemple, l'estimation de la surface de rectangles) entraînerait une amélioration globale des
capacités d'observation ou de raisonnement sur d'autres formes géométriques ou d'autres tâches
intellectuelles.\cite{II-1-2-ref6} Leurs résultats furent sans appel et constituèrent un séisme
pédagogique : l'amélioration constatée était strictement limitée à la tâche entraînée. Les sujets
devenaient excellents pour estimer des rectangles, mais ne montraient aucune amélioration
significative pour estimer des cercles ou des triangles, et encore moins pour des tâches de logique
verbale.\cite{II-1-2-ref14}

Thorndike a ainsi formulé la « théorie des éléments identiques » (identical elements theory), qui
stipule que le transfert d'apprentissage d'une situation A vers une situation B ne se produit que
dans la mesure où A et B partagent des éléments concrets communs.\cite{II-1-2-ref6} Autrement dit,
étudier la grammaire latine aide à apprendre la grammaire espagnole (car elles partagent des racines
lexicales et des structures syntaxiques communes), mais n'améliore pas intrinsèquement la capacité à
résoudre des équations différentielles, à jouer aux échecs ou à structurer un raisonnement
juridique, car ces activités ne partagent que peu ou pas d'éléments cognitifs identiques avec la
morphologie latine.\cite{II-1-2-ref16} L'idée que le grec ancien agirait comme un aiguiseur
universel de l'intelligence s'est révélée, dès 1901, être une illusion cognitive dépourvue de
fondement expérimental.



\subsection{La science cognitive moderne : la distinction transfert proche vs transfert lointain}

Les recherches contemporaines en neurosciences et en psychologie cognitive ont confirmé et affiné
les conclusions de Thorndike, en introduisant une distinction conceptuelle cruciale pour notre sujet
: celle entre le « transfert proche » (\textit{near transfer}) et le « transfert lointain »
(\textit{far transfer}).\cite{II-1-2-ref7}

\begin{itemize}\item \textbf{Le Transfert Proche :} Il se caractérise par l'application de compétences acquises dans un contexte d'apprentissage à un contexte très similaire. Par exemple, l'analyse syntaxique d'une phrase complexe en grec (identification du sujet, du verbe, des compléments) favorise effectivement le développement de la conscience métalinguistique et aide à l'analyse grammaticale en français ou à l'apprentissage de l'allemand.\cite{II-1-2-ref17} C'est un transfert réel, documenté, mais strictement confiné au domaine linguistique.\item \textbf{Le Transfert Lointain :} Il désigne la capacité hypothétique d'appliquer une compétence apprise dans un domaine spécifique (ex: la morphologie verbale grecque) à un domaine totalement différent et structurellement distinct (ex: la logique mathématique, la planification stratégique, ou l'intelligence fluide générale). La littérature scientifique actuelle est quasi-unanime : le transfert lointain est extrêmement rare, difficile à provoquer, et souvent inexistant dans les conditions éducatives standard.\cite{II-1-2-ref7}\end{itemize}Les méta-analyses récentes portant sur l'entraînement cognitif (\textit{brain training}), la
pratique des échecs ou l'apprentissage de la musique aboutissent aux mêmes conclusions décevantes
pour les partisans de la discipline formelle. Jouer aux échecs améliore la capacité à jouer aux
échecs, mais ne rend pas les élèves significativement plus intelligents en mathématiques ou en
lecture une fois les facteurs confondants contrôlés.\cite{II-1-2-ref8} De même, les études sur la
mémoire de travail montrent que l'entraînement sur une tâche spécifique (comme le \textit{n-back}
task) n'augmente pas l'intelligence fluide générale (\textit{Gf}).\cite{II-1-2-ref21}

Appliqué au cas du grec ancien, cela signifie que la gymnastique mentale consistant à jongler avec
les cas, les modes, les aspects verbaux et les dialectes développe des compétences hautement
spécialisées de traitement du langage, de mémoire morphologique et d'attention au détail textuel.
Cependant, elle ne crée pas une « compétence de rigueur » générique qui s'exporterait magiquement
vers la géométrie non-euclidienne ou l'algorithmie.\cite{II-1-2-ref22} Le cerveau ne fonctionne pas
comme un muscle unique, mais plutôt comme un ensemble de réseaux spécialisés dont les connexions ne
se renforcent que par la sollicitation spécifique des circuits concernés.



\subsection{La spécificité de domaine (\textit{domain specificity}) et l'architecture cognitive}

Un concept clé pour comprendre l'absence de ce transfert espéré est la « spécificité de domaine »
(\textit{domain specificity}).\cite{II-1-2-ref23} La cognition humaine tend à être compartimentée en
modules fonctionnels qui ont évolué pour traiter des types d'informations spécifiques.

La théorie de la modularité de l'esprit, bien que débattue dans ses versions les plus strictes
(Fodor), trouve un écho dans la distinction entre les systèmes de mémoire déclarative et
procédurale, telle que modélisée par Ullman dans le modèle DP (\textit{Declarative/Procedural
model}).\cite{II-1-2-ref25} Selon ce modèle :

\begin{itemize}\item Le \textbf{lexique mental} (les mots, les irrégularités) dépend de la mémoire déclarative (lobe temporal), utilisée pour stocker des faits et des événements.\item La \textbf{grammaire mentale} (les règles combinatoires) dépend de la mémoire procédurale (ganglions de la base, cortex frontal), utilisée pour apprendre des séquences motrices et cognitives.\cite{II-1-2-ref25}\end{itemize}Bien que les mathématiques sollicitent également des processus procéduraux, les réseaux neuronaux
impliqués dans le traitement syntaxique du langage (aire de Broca, gyrus frontal inférieur) sont
distincts de ceux impliqués dans le raisonnement numérique et la manipulation algébrique (sillon
intrapariétal, cortex préfrontal dorsolatéral).\cite{II-1-2-ref28} Des études récentes en
neuroimagerie (IRMf) montrent que le réseau du langage et le réseau de la logique mathématique sont
largement dissociés dans le cerveau adulte : il est possible d'avoir des déficits massifs dans l'un
tout en préservant l'autre (comme chez certains aphasiques capables de raisonnement logique
complexe, ou des dyscalculiques parfaitement éloquents).\cite{II-1-2-ref28}

Croire que l'on peut entraîner la « logique mathématique » en faisant de la « logique grammaticale »
revient donc à une erreur de catégorie neurobiologique. La « rigueur » du latiniste est une rigueur
philologique (respect des règles d'accord, précision lexicale, concordance des temps), tandis que la
« rigueur » mathématique repose sur l'abstraction, la déduction axiomatique et la visualisation
spatiale.\cite{II-1-2-ref29} Bien que le langage courant utilise le même mot (« rigueur ») pour
désigner ces deux qualités, créant une confusion sémantique, les processus cognitifs sous-jacents ne
se chevauchent que marginalement.\cite{II-1-2-ref31}

Si la théorie cognitive prédit une absence de transfert, comment expliquer le constat sociologique
récurrent selon lequel les élèves latinistes et hellénistes obtiennent de meilleurs résultats
scolaires globaux, et particulièrement en mathématiques? C'est ici qu'intervient l'analyse critique
des données empiriques, qui permet de distinguer la causalité pédagogique de la corrélation sociale.



\subsection{L'étude de nuremberg (haag \& stern, 2003\) : une réfutation contrôlée}

L'une des études les plus rigoureuses et les plus citées dans ce domaine a été menée par Ludwig Haag
et Elsbeth Stern à l'Université de Nuremberg. Dans un contexte éducatif allemand où le choix du
latin reste un marqueur fort, ils ont cherché à vérifier empiriquement si l'apprentissage du latin
améliorait réellement les capacités de raisonnement logique (déductif et inductif) par rapport à
l'apprentissage d'une langue vivante ou à l'absence de latin.\cite{II-1-2-ref16}

Leurs conclusions sont dévastatrices pour la thèse utilitariste du « muscle mental » :

1. \textbf{Aucun gain de QI attribuable au latin :} Les chercheurs n'ont trouvé aucune différence
significative dans l'évolution du QI verbal ou non-verbal entre les élèves ayant étudié le latin et
ceux ne l'ayant pas étudié, une fois les différences initiales d'intelligence contrôlées. Le latin
ne rend pas « plus intelligent » ; ce sont les élèves intelligents qui choisissent le
latin.\cite{II-1-2-ref22}

2. \textbf{Absence de transfert logique :} Sur les tests de raisonnement déductif et inductif
(proches de la logique mathématique), les latinistes ne surperforment pas leurs pairs non-latinistes
après plusieurs années d'étude intensive. La pratique de la version latine n'a pas développé de
compétence logique généraliste.\cite{II-1-2-ref16}

3. \textbf{Désavantage linguistique paradoxal :} Plus surprenant encore, l'étude a suggéré que pour
l'apprentissage ultérieur de l'espagnol, les élèves ayant étudié le français ou l'anglais s'en
sortaient mieux que les latinistes. Ce résultat contredit frontalement l'idée reçue selon laquelle
le latin serait la « base logique » indispensable à l'apprentissage des langues romanes. Les
interférences cognitives et la méthode d'enseignement trop analytique du latin (grammaire-
traduction) semblaient même handicaper l'acquisition communicative d'une nouvelle
langue.\cite{II-1-2-ref33}

Haag et Stern concluent que les bénéfices cognitifs supposés du latin relèvent du mythe pédagogique.
La supériorité apparente des latinistes dans les statistiques brutes ne provient pas du contenu de
l'enseignement du latin, mais des caractéristiques préexistantes de la population qui choisit cette
option.



\subsection{Les travaux longitudinaux d'alexandra vereeck : la preuve par la sélection}

Plus récemment, les travaux longitudinaux menés par Alexandra Vereeck en Flandre (Belgique) ont
apporté des éclairages définitifs sur le mécanisme de cette illusion statistique.\cite{II-1-2-ref9}
Le système éducatif flamand, très proche du système français dans sa valorisation des filières
classiques, offre un terrain d'étude idéal.

L'étude de Vereeck montre que les élèves qui choisissent le latin présentent, \textit{dès le début
du secondaire} (avant même d'avoir été exposés significativement à la langue), des scores
significativement plus élevés en intelligence générale, en compétences linguistiques natives et en
conscience métalinguistique.\cite{II-1-2-ref9} Il s'agit d'un phénomène massif de
\textbf{présélection} (\textit{self-selection bias}). Les élèves performants, souvent issus de
milieux socio-économiques favorisés et encouragés par des parents informés, s'orientent vers les
filières classiques parce qu'elles sont réputées difficiles et prestigieuses.

L'apport crucial de l'étude longitudinale est de montrer que l'écart de performance entre les
latinistes et les non-latinistes, bien que présent en début de cursus, ne s'accroît pas de manière
disproportionnée au fil des ans. Si le latin agissait comme un « accélérateur cognitif » ou un
entraînement à la rigueur, on devrait observer une divergence croissante des courbes de performance
en faveur des latinistes. Or, pour la plupart des compétences cognitives et mathématiques mesurées,
l'écart reste stable ou évolue parallèlement à celui du groupe témoin à intelligence initiale
égale.\cite{II-1-2-ref9} L'avantage est donc préexistant à l'enseignement, et non causé par celui-
ci.



\subsection{Le facteur socio-économique et l'habitus : la variable cachée}

L'analyse des résultats scolaires ne peut faire l'économie de la variable sociologique. Les études
montrent systématiquement que lorsque l'on contrôle rigoureusement le statut socio-économique (SSE)
des parents et le capital culturel familial, l'avantage résiduel des étudiants en langues anciennes
diminue drastiquement, voire disparaît.\cite{II-1-2-ref37}

Les élèves de grec et de latin bénéficient majoritairement d'un environnement familial riche en
capital culturel, propice au développement de l'abstraction, du vocabulaire soutenu et des
dispositions scolaires. C'est cet environnement, combiné à une scolarisation dans des classes
souvent plus calmes, plus motivées et bénéficiant de meilleurs enseignants (effet de pair et effet
maître), qui favorise la réussite en mathématiques.\cite{II-1-2-ref10} Pierre Bourdieu parlait d'un
« habitus » scolaire ; les humanités classiques fonctionnent comme un marqueur et un amplificateur
de cet habitus, attirant ceux qui possèdent déjà les codes de l'excellence
scolaire.\cite{II-1-2-ref39} Attribuer leur réussite en mathématiques à la grammaire grecque revient
à confondre l'étiquette du flacon avec le contenu du breuvage.



\subsection{Tableau 1 : synthèse des études empiriques sur le transfert cognitif du latin/grec}

\begin{table}[h!]\small\centering\begin{tabular}{| p{2cm} | p{2.2cm} | p{2.2cm} | p{2.2cm} |}\hline\textbf{Étude} & \textbf{Méthodologie} & \textbf{Résultats Principaux} & \textbf{Conclusion sur le Transfert} \\ \hline\textbf{Thorndike \& Woodworth (1901)} & Expérimentale (estimation de surfaces) & Amélioration limitée à la tâche spécifique. Pas de généralisation. & \textbf{Réfutation} de la discipline formelle. \\ \hline\textbf{Haag \& Stern (2003)} & Longitudinale (Allemagne) & Aucun gain de QI ou de logique déductive chez les latinistes vs non-latinistes. & \textbf{Absence} de transfert cognitif vers la logique. \\ \hline\textbf{Vereeck (2020-2023)} & Longitudinale (Belgique, N=1731) & Avantage cognitif présent \textit{avant} l'apprentissage (biais de sélection). Pas d'accroissement de l'écart. & \textbf{Illusion} statistique due à la présélection. \\ \hline\textbf{Méta-analyses (Sala et al., 2019)} & Méta-analyse (Entraînement cognitif) & Transfert lointain quasi-nul pour les jeux d'échecs, musique, mémoire de travail. & \textbf{Scepticisme} sur tout transfert lointain. \\ \hline\end{tabular}\end{table}L'argument central de la justification utilitariste repose sur l'analogie : la grammaire ancienne
serait une « logique » en soi, une algèbre linguistique qui prépare structurellement à l'algèbre
mathématique. Cette analogie, séduisante en surface, ne résiste pas à l'analyse épistémologique,
linguistique et cognitive approfondie.



\subsection{Nature de la logique : formelle vs naturelle}

Il existe une différence de nature, et non simplement de degré, entre la logique mathématique
(axiomatique) et la « logique » grammaticale d'une langue naturelle.

\begin{itemize}\item \textbf{La Logique Mathématique} est un système formel hypothético-déductif. Elle opère sur des symboles univoques définis par des axiomes stricts. Dans un système formel, la non-contradiction est absolue (A ne peut être non-A). La vérité d'une conclusion dépend exclusivement de la validité de la chaîne de déduction, indépendamment du contexte extérieur.\cite{II-1-2-ref40}\item \textbf{La Logique Grammaticale}, même dans une langue aussi structurée que le grec ancien, est fondamentalement différente. C'est un système descriptif et conventionnel, issu de l'usage et de l'évolution historique, et non d'une création axiomatique \textit{ex nihilo}. Elle est saturée d'exceptions, d'irrégularités, de nuances sémantiques contextuelles et d'ambiguïtés.\cite{II-1-2-ref42} La « règle » grammaticale n'est pas une loi universelle comme la gravité ou le théorème de Pythagore ; elle est une généralisation statistique de l'usage avec une marge de flou incompressible.\end{itemize}Comme le soulignent les chercheurs en Intelligence Artificielle et en Traitement du Langage Naturel
(NLP), traduire le langage naturel en logique formelle est une tâche d'une complexité inouïe
précisément parce que le langage n'obéit pas à la rigueur stricte des
mathématiques.\cite{II-1-2-ref43} Les tentatives de modéliser la grammaire comme un calcul logique
pur (type grammaires de Montague) se heurtent toujours à la plasticité du sens et à la pragmatique.
Faire de l'analyse grammaticale, c'est donc exercer une activité de \textbf{classification} et de
\textbf{reconnaissance de motifs} (\textit{pattern matching}) dans un système bruité, ce qui est
cognitivement très différent de la \textbf{démonstration} mathématique dans un système clos et
parfait.\cite{II-1-2-ref30}



\subsection{La théorie de la charge cognitive (\textit{cognitive load theory}) et l'échec de la traduction}

L'enseignement traditionnel du grec (méthode grammaire-traduction) est souvent défendu pour sa
difficulté même : « ça fait mal à la tête, donc ça travaille ». Or, la Théorie de la Charge
Cognitive, développée par John Sweller, suggère que cette difficulté peut être artificielle et
contre-productive.\cite{II-1-2-ref45}

Lorsqu'un élève traduit une phrase grecque complexe (ex: du Platon ou du Thucydide) en se focalisant
intensément sur l'identification analytique des cas (nominatif, génitif, datif), des temps et des
modes, sa mémoire de travail est saturée par des tâches de « bas niveau » (décodage morphologique).
Cette surcharge cognitive (\textit{extraneous cognitive load}) consomme les ressources
attentionnelles qui devraient être allouées à la compréhension du sens profond, de la structure
argumentative ou de la beauté littéraire du texte.\cite{II-1-2-ref47}

Au lieu de développer une pensée fluide et rigoureuse, l'élève s'épuise dans un traitement de
données mécanique et séquentiel qui s'apparente plus à du déchiffrage laborieux qu'à du raisonnement
de haut vol.\cite{II-1-2-ref15} Loin d'être une préparation à la fluidité mathématique, cette
méthode entraîne le cerveau à une forme de lenteur analytique et de vérification constante qui n'est
pas celle requise pour l'intuition mathématique ou la résolution de problèmes dynamiques, où la
fluidité et l'automatisation des procédures de base sont essentielles.\cite{II-1-2-ref48}



\subsection{L'impasse pédagogique : de la grammaire-traduction aux méthodes vivantes}

L'argument de la rigueur a servi historiquement à justifier le maintien de la méthode « Grammaire-
Traduction » contre les méthodes plus communicatives, immersives ou intuitives (comme la méthode
Polis, le TPR \- \textit{Total Physical Response} ou les approches basées sur l'input compréhensible
de Krashen).\cite{II-1-2-ref50}

Les recherches en acquisition des langues secondes (SLA) montrent pourtant que la connaissance
explicite et métalinguistique des règles grammaticales (le « Moniteur » dans la théorie de Krashen)
n'est pas nécessaire pour comprendre ou produire une langue, et qu'elle peut même entraver la
fluidité en imposant un contrôle conscient trop lourd.\cite{II-1-2-ref53} En insistant sur l'analyse
(le « pourquoi » c'est un accusatif) au détriment de l'exposition massive à la langue (le « input »
compréhensible de Paul Nation, qui requiert 98\% de vocabulaire connu pour une lecture fluide), on
forme des élèves capables de disséquer une phrase morte comme un cadavre, mais incapables de penser
\textit{dans} la langue.\cite{II-1-2-ref55}

Si l'objectif réel était la rigueur logique, l'enseignement de la logique formelle pure, des
statistiques ou de la programmation informatique serait infiniment plus efficace, direct et
pertinent que le détour laborieux par la déclinaison thématique du grec ancien.\cite{II-1-2-ref49}
Persister à utiliser le grec comme un ersatz de cours de logique est non seulement une erreur
pédagogique, mais un détournement de la finalité culturelle de la discipline.

Si les preuves scientifiques du transfert sont si faibles et si l'analogie logique est boiteuse,
pourquoi cet argument perdure-t-il avec tant de vigueur dans le discours éducatif et parental? La
réponse ne se trouve pas dans la neurologie, mais dans l'histoire politique de l'éducation et la
sociologie des institutions.



\subsection{La crise de 1880 et le « sauvetage » des humanités par la gymnastique}

À la fin du XIXe siècle, en France, l'enseignement secondaire classique entre en crise
existentielle. Sous la IIIe République, des réformateurs républicains comme Jules Ferry et des
intellectuels comme le linguiste Michel Bréal remettent en cause l'hégémonie absolue du latin et du
grec, jugés inadaptés à la modernité industrielle, scientifique et démocratique.\cite{II-1-2-ref12}
Des exercices emblématiques comme le « vers latin » ou le discours latin sont supprimés ou
marginalisés.

C'est à ce moment précis de fragilité institutionnelle que l'argumentaire bascule. Puisque l'utilité
directe (écrire en latin pour l'Église, diplomatie, sciences) disparaît, il faut inventer une
nouvelle utilité, plus abstraite et inattaquable. Michel Bréal lui-même, bien que critique des
méthodes anciennes, et d'autres défenseurs des humanités, vont théoriser la « gymnastique de
l'esprit ».\cite{II-1-2-ref3} Le latin et le grec deviennent des disciplines « gratuites » et «
désintéressées », dont la seule fin est la formation du caractère et de l'intelligence formelle.

André Chervel, dans son \textit{Histoire des disciplines scolaires} et ses travaux sur la grammaire,
a magistralement montré comment la grammaire scolaire a été instrumentalisée pour donner une
consistance « disciplinaire » et sérieuse à l'enseignement du français et des langues anciennes,
indépendamment de l'usage réel de la langue.\cite{II-1-2-ref11} On a maintenu des exercices
artificiels (comme le thème d'imitation ou la version ultra-analytique) non pas parce qu'ils étaient
efficaces pour apprendre la langue, mais parce qu'ils étaient difficiles, classants, et permettaient
de hiérarchiser les élèves selon une échelle de « mérite » intellectuel abstrait.\cite{II-1-2-ref62}
La difficulté est devenue une fin en soi, garante de la valeur formatrice.



\subsection{La distinction et la prophétie autoréalisatrice (bourdieu)}

Pierre Bourdieu, dans ses œuvres majeures \textit{La Distinction} et \textit{La Reproduction}, a
offert la clé sociologique de ce phénomène. L'enseignement des langues anciennes fonctionne comme un
mécanisme de distinction sociale raffiné. Le « culte de la rigueur » inutile est un marqueur de
classe : c'est précisément parce que c'est « gratuit » (sans utilité économique immédiate,
contrairement à l'anglais ou à la comptabilité) que c'est prestigieux. Cela signale une distance au
besoin, un luxe temporel et intellectuel réservé à l'élite.\cite{II-1-2-ref39}

Cette dynamique crée une prophétie autoréalisatrice parfaite, que le sociologue Eric Mangez décrit
également dans le contexte belge 10 :

1. Le système scolaire et les familles étiquettent les élèves les plus performants (souvent issus de
milieux favorisés) et les orientent massivement vers les filières classiques (Latin-Grec ou Latin-
Maths).

2. On inculque à ces élèves l'idée que le grec va former leur esprit à une rigueur supérieure.

3. Ces élèves réussissent brillamment en mathématiques et dans les études supérieures (grâce à leur
niveau initial, leur capital culturel et leur confiance en eux).

4. Le système éducatif et les parents concluent a posteriori : « Vous voyez? Le grec a formé leur
esprit mathématique\! »

En réalité, le grec n'a pas \textit{causé} la réussite mathématique ; il a servi de \textit{signe de
reconnaissance} et de \textit{chambre de tri} pour une élite scolaire préexistante. L'argument
utilitariste (« ça sert aux maths ») est un masque méritocratique posé sur une fonction de
reproduction sociale. Il permet de justifier les inégalités de parcours par une supposée inégalité
de « formation de l'esprit », naturalisant ainsi les avantages sociaux.



\subsection{L'excellence pour tous : une nouvelle rhétorique pour un vieux mécanisme?}

Aujourd'hui, face à la démocratisation de l'enseignement, l'argumentaire a légèrement glissé vers
l'idée d'« excellence pour tous » ou de remédiation linguistique (le latin pour mieux comprendre
l'orthographe française). Si le transfert lexical (étymologie, vocabulaire technique) est mieux
attesté (transfert proche) et pertinent 17, l'argument de la « rigueur mathématique » reste brandi
par les parents « initiés » et certains enseignants pour justifier le maintien des options
classiques dans les lycées d'élite ou pour éviter les collèges de secteur jugés moins
bons.\cite{II-1-2-ref65}

L'existence de filières comme la « Filière C » (Maths) et la « Filière A » (Lettres) en France a
longtemps cristallisé cette hiérarchie, où le latin servait de ticket d'entrée officieux pour
l'élite scientifique.\cite{II-1-2-ref66} Même après les réformes, l'ombre portée de cette hiérarchie
subsiste : faire du grec reste un signal fort d'appartenance au groupe des « bons élèves » destinés
aux classes préparatoires, indépendamment de tout amour pour la culture antique. C'est une stratégie
de contournement de la carte scolaire et de maintien de l'entre-soi, habillée de justifications
cognitives que la science dément.

L'analyse approfondie et croisée des données cognitives, historiques et sociologiques mène à une
conclusion univoque : il n'y a pas de transfert automatique de « rigueur » entre l'apprentissage
scolaire du grec ancien et les mathématiques.

1. \textbf{L'échec scientifique de la théorie formelle :} Le cerveau n'est pas un muscle unique ;
les compétences logiques sont largement spécifiques au domaine et ne se transfèrent pas magiquement
de la grammaire à l'algèbre.

2. \textbf{L'illusion statistique de la réussite :} La corrélation observée entre grec et maths est
le fruit d'une sélection sociale et cognitive en amont (biais de sélection), et non d'une causalité
pédagogique interne à la discipline.

3. \textbf{L'incompatibilité épistémologique :} La rigueur grammaticale (classificatoire, normative,
sémantique) diffère radicalement dans sa nature de la rigueur mathématique (axiomatique, déductive,
formelle).

4. \textbf{L'origine défensive et sociale de l'argument :} La justification utilitariste est une
construction historique de la fin du XIXe siècle destinée à sauver une discipline dont l'utilité
sociale directe s'effondrait, en la transformant en outil de distinction des élites.

Affirmer aujourd'hui que « faire du grec rend rigoureux en maths » est non seulement une erreur
scientifique, mais aussi une impasse stratégique pour les défenseurs des humanités. En promettant
des « super-pouvoirs » cognitifs que la discipline ne peut pas délivrer, on l'expose à la déception
et à la critique légitime des sciences de l'éducation. De plus, cette vision instrumentalise les
textes antiques, les réduisant à des prétextes pour des exercices de logique douteux, au lieu de les
aborder pour ce qu'ils sont : des œuvres de pensée, de beauté et d'humanité.

Le grec ancien et le latin doivent être enseignés et défendus pour leur valeur intrinsèque : pour
comprendre l'histoire de notre langue, pour accéder directement à des textes fondateurs (Homère,
Platon, Sophocle, les Évangiles), pour développer une conscience historique et esthétique, et pour
le plaisir intellectuel de déchiffrer une altérité culturelle. Ils n'ont pas besoin d'être les
servantes des mathématiques pour être légitimes. Libérer l'enseignement des langues anciennes de ce
mythe utilitariste, c'est peut-être enfin leur permettre d'être enseignées comme des langues
vivantes de culture, accessibles et passionnantes, plutôt que comme des instruments de torture
sélective.

